\appendix

\section{Supplementary Material}
\label{sec:appendix}

\subsection{CKA Analysis Method}
\label{sec:cka_appendix}

We employ Centered Kernel Alignment (CKA)~\cite{kornblith2019} alongside RSA as a complementary perspective.
While RSA captures local geometric alignment via rank-order preservation, CKA measures global representational similarity through kernel-based comparisons.

We compute both \emph{linear CKA}, which uses a linear kernel and is sensitive to second-order statistics, and \emph{RDM kernel CKA}, which operates on the same distance matrices as RSA but aggregates them differently.
Using both variants tests whether conclusions about representation structure are robust to the choice of similarity metric.

Linear CKA between high-dimensional representations (4{,}096-D for LLM layers) and the 7-D action space is subject to a known dimension mismatch: the linear kernel in the low-dimensional space may not capture the same structure as in the high-dimensional space.
RDM kernel CKA avoids this issue by operating on distance matrices of equal size regardless of the original feature dimensionality.

\subsection{Metric Sensitivity: RSA versus CKA}
\label{sec:metric_sensitivity_appendix}

Linear CKA~\cite{kornblith2019} reverses the RSA ordering: untrained LLM ($\text{CKA} = 0.089$) $>$ trained LLM ($0.083$) $>$ DinoSigLIP+Projector ($0.081$) $>$ DINOv2 ($0.066$), partially inverting the self-supervised advantage (Figure~\ref{fig:metric_comparison}).
RDM kernel CKA partially resolves this discrepancy: trained ($0.260$) $\approx$ Projector ($0.257$) $>$ untrained ($0.246$) $>$ DINOv2 ($0.241$).
This divergence reflects the known sensitivity of representation similarity metrics to different geometric properties.

At the per-layer level, linear CKA between image\_mean and the action space rises monotonically from L0 ($0.081$) to L32 ($0.098$), showing no inverted-U pattern.
RSA captures rank-order geometric alignment; CKA captures global linear structure.
The dissociation indicates that the geometric refinement in early layers operates on the relative arrangement of representations (affecting RSA) without altering the overall covariance structure (captured by CKA).

\begin{figure}[t]
\centering
\includegraphics[width=\columnwidth]{figures/fig6.pdf}
\caption{Three similarity metrics across representation sources. \textbf{(a)}~RSA ranks DINOv2 cls ($0.692$) $>$ Fused ($0.634$) $>$ Projector ($0.615$). \textbf{(b)}~Linear CKA reverses the ordering: Untrained ($0.089$) $>$ Trained ($0.083$) $>$ Projector ($0.081$) $>$ DINOv2 ($0.066$). \textbf{(c)}~RDM kernel CKA: Trained ($0.260$) $\approx$ Projector ($0.257$) $>$ Untrained ($0.246$) $>$ DINOv2 ($0.241$). The metric divergence underscores that RSA and CKA capture different aspects of representational structure.}
\label{fig:metric_comparison}
\end{figure}

\subsection{Branch-Split Detailed Analysis}
\label{sec:branch_split_appendix}

To disentangle the contributions of the two encoder streams within DinoSigLIP, we extract representations from the DINOv2 branch (1{,}024-D) and SigLIP branch (1{,}152-D) separately before fusion.
The two branches yield near-identical RSA: DINOv2 branch $\rho = 0.619$, SigLIP branch $\rho = 0.621$.
Fusion produces a representation ($\rho = 0.634$) that exceeds either branch alone, indicating that concatenation captures complementary structure.
The MLP projector then compresses this signal: DinoSigLIP+Projector achieves $\rho = 0.615$, a drop of $\Delta = 0.019$ from the fused 2{,}176-D input.

\subsection{Additional RSA Values}
\label{sec:additional_rsa}

DINOv2 \texttt{patch\_mean} achieves RSA $\rho = 0.631$, between the \texttt{cls\_token} ($0.692$) and the fused DinoSigLIP ($0.634$).
VLA untrained readout peaks at L1 ($\rho = 0.626$).
Instruction invariance holds across all conditions (Figure~\ref{fig:factorial}, panel b): maximum $|\Delta| < 0.006$.

\subsection{Cross-Benchmark Validation: LIBERO-Spatial}
\label{sec:cross_benchmark_appendix}

To assess generalization beyond in-distribution tasks, we repeat the core RSA and probe analyses on LIBERO-Spatial (10 out-of-distribution tasks requiring spatial reasoning).
Figure~\ref{fig:cross_benchmark} shows per-layer RSA and $R^2$ profiles for both LIBERO-Goal and LIBERO-Spatial side by side.

The inverted-U profile replicates on LIBERO-Spatial: trained image\_mean RSA peaks at L14 ($\rho = 0.691$) and $R^2$ at L13 ($0.696$), with peaks shifting ${\sim}2$ layers deeper than LIBERO-Goal (RSA L12, $R^2$ L8--L10)---consistent with increased task difficulty requiring deeper processing.
Training produces a significant geometric gain (best-vs-best $\Delta = +0.023$, 95\% CI $[0.012, 0.036]$, $p < 0.001$), larger than the LIBERO-Goal effect ($\Delta = +0.017$, CI includes zero) and confirming that the n.s.\ result on Goal reflects limited statistical power.

Two notable differences emerge.
First, the absolute RSA values are higher on Spatial ($0.691$ vs.\ $0.628$), possibly reflecting greater action diversity in spatial tasks that produces a more discriminative action RDM.
Second, RSA--$R^2$ peak convergence does not replicate: readout RSA peaks at L15 while $R^2$ peaks at L21 (6-layer gap), contrasting with the near-alignment on Goal ($|\Delta_\text{layer}| = 1$).
This suggests that peak convergence is benchmark-specific rather than a general consequence of training.

\begin{figure}[t]
\centering
\includegraphics[width=\columnwidth]{figures/fig7.pdf}
\caption{Cross-benchmark validation: per-layer RSA and $R^2$ profiles on LIBERO-Goal (left) vs.\ LIBERO-Spatial (right). The inverted-U replicates on both benchmarks. LIBERO-Spatial shows a significant training effect (RSA best-vs-best $\Delta = +0.023$, $p < 0.001$) and deeper peak layers (RSA L14, $R^2$ L13), but RSA--$R^2$ peak convergence does not replicate (6-layer gap).}
\label{fig:cross_benchmark}
\end{figure}

\subsection{Robustness Analyses}
\label{sec:robustness_appendix}

\subsubsection{Distance Metric Invariance}
\label{sec:metric_invariance}

The inverted-U profile is robust to distance metric choice: cosine action distance yields identical peak layers (L12) and Pearson correlation $r > 0.989$ with the Euclidean profile reported in the main text (Table~\ref{tab:metric_invariance}).

\begin{table}[h]
\centering
\caption{Peak layer comparison between Euclidean and cosine action distances. Profile $r$ denotes Pearson correlation across all 32 layers.}
\label{tab:metric_invariance}
\small
\begin{tabular}{lccc}
\toprule
Condition & Eucl.\ Peak & Cos.\ Peak & Profile $r$ \\
\midrule
Trained image\_mean & L12 & L12 & 0.990 \\
Trained last\_preaction & L12 & L13 & 0.990 \\
Untrained image\_mean & L0 & L0 & 0.995 \\
Untrained last\_preaction & L1 & L1 & 0.992 \\
\bottomrule
\end{tabular}
\end{table}

\subsubsection{Per-DOF Decomposition}
\label{sec:perdof}

RSA decomposition by action dimension confirms the inverted-U is not driven by any single degree of freedom: xyz translation (peak L13, $\rho = 0.193$), rotation (peak L12, $\rho = 0.175$), and gripper (peak L13, $\rho = 0.106$) all exhibit non-monotonic profiles.
Per-DOF RSA values are lower than the 7-D RSA because single-dimension distances carry less discriminative information than the full action space.

\subsubsection{Permutation Test}
\label{sec:permutation}

Permutation testing (1{,}000 shuffles) confirms all observed RSA values are highly significant: observed $\rho$ ($0.42$--$0.63$) exceeds the null distribution maximum ($0.08$--$0.10$) by $15$--$23$ standard deviations (all $p < 0.001$).

\subsubsection{Probe Seed Sensitivity}
\label{sec:seed_sensitivity}

Action probe $R^2$ is highly stable across random seeds.
Over 5 seeds (trained image\_mean, full-layer scan), per-layer $R^2$ standard deviation is uniformly below $0.01$ (mean $0.004$), and the peak $R^2$ varies by less than $0.4\%$ ($0.655 \pm 0.002$; Table~\ref{tab:seed_sensitivity}).
Peak layer concentrates in the L8--L10 region ($4/5$ seeds), confirming the inverted-U shape is not a seed artifact.
Seed 1 peaks at L16, but with $R^2$ difference ${<}0.002$---the inverted-U plateau (L8--L16) is sufficiently flat that argmax is sensitive to minor noise while the profile shape is invariant.

\begin{table}[h]
\centering
\caption{Probe $R^2$ peak across 5 random seeds (trained image\_mean).}
\label{tab:seed_sensitivity}
\small
\begin{tabular}{lcc}
\toprule
Seed & Peak Layer & Peak $R^2$ \\
\midrule
0 & L9 & 0.654 \\
1 & L16 & 0.656 \\
2 & L9 & 0.653 \\
3 & L10 & 0.655 \\
42 & L9 & 0.655 \\
\midrule
Mean$\pm$std & $10.6\pm2.7$ & $0.655\pm0.002$ \\
\bottomrule
\end{tabular}
\end{table}

\subsubsection{Demonstration Subsample Sensitivity}
\label{sec:demo_sensitivity}

To verify that the inverted-U pattern is not an artifact of sample size, we subsample 25 from 50 demonstrations per task and recompute both RSA and probe profiles (Table~\ref{tab:demo_sensitivity}).
The RSA peak layer remains at L12 with a profile correlation of $r = 0.970$ between 50- and 25-demo conditions, confirming high stability.
The probe peak shifts from L9 to L11 (within the L8--L16 plateau) with $r = 0.819$, reflecting greater sensitivity of $R^2$ to training set size.
Both metrics show reduced absolute values (RSA $-3.6\%$, $R^2$ $-7.6\%$) as expected from halved data, but the non-monotonic profile shape is preserved.

\begin{table}[h]
\centering
\caption{Inverted-U stability under demonstration subsampling (trained image\_mean).}
\label{tab:demo_sensitivity}
\small
\begin{tabular}{lccc}
\toprule
Metric & 50 demos & 25 demos & $\Delta$ \\
\midrule
RSA peak layer & L12 & L12 & 0 \\
RSA peak $\rho$ & 0.628 & 0.606 & $-3.6\%$ \\
$R^2$ peak layer & L9 & L11 & $+2$ \\
$R^2$ peak value & 0.659 & 0.609 & $-7.6\%$ \\
RSA profile $r$ & --- & 0.970 & --- \\
$R^2$ profile $r$ & --- & 0.819 & --- \\
\bottomrule
\end{tabular}
\end{table}

\subsection{Training Effect Size Summary}
\label{sec:training_effect_summary}

Table~\ref{tab:training_effect_summary} consolidates training effects across all metrics.
Training effects on image\_mean metrics are modest (RSA $\Delta = +0.017$, $R^2$ $\Delta = +0.015$), but effects on last\_preaction metrics are substantial (routing $+18$pp, attention $+46$pp, $R^2$ $\Delta = +0.069$).
This asymmetry confirms that training primarily reshapes how the readout token integrates information rather than altering the visual representations carried by image tokens.

\begin{table}[h]
\centering
\caption{Training effect sizes across metrics. Image\_mean effects are modest; last\_preaction effects are large, confirming that training targets readout-token information integration.}
\label{tab:training_effect_summary}
\small
\begin{tabular}{lccc}
\toprule
Metric & Trained & Untrained & $\Delta$ \\
\midrule
RSA peak (image\_mean) & 0.628 (L12) & 0.612 (L0) & $+$0.017 \\
$R^2$ peak (image\_mean) & 0.655 (L8) & 0.640 (L0) & $+$0.015 \\
$R^2$ peak (last\_preaction) & 0.678 (L13) & 0.609 (L32) & $+$0.069 \\
Instruction routing L31 & 100\% & 82\% & $+$18pp \\
Instruction attention L1 & 58\% & 12\% & $+$46pp \\
Progress classification & 86.3\% (L20) & 85.3\% (L4) & $+$1pp \\
\bottomrule
\end{tabular}
\end{table}

\subsection{Instruction Routing Probes}
\label{sec:routing_appendix}

\subsubsection{Cross-Instruction Classification: Dense Layer Profile}
\label{sec:dense_routing}

Table~\ref{tab:dense_routing} shows the dense per-layer profile of cross-instruction classification (10-way) at the readout token for layers L0--L8.
The step-function onset between L0 ($10\%$, chance) and L1 (${\geq}99.6\%$) reflects the high discriminability of instruction embeddings from the frozen tokenizer: even incidental attention (${\sim}12\%$) suffices for linear classification at L1, though only trained deliberate routing (58\% attention) sustains the signal through all 31 layers.

\begin{table}[h]
\centering
\caption{Cross-instruction 10-way classification accuracy at readout token, L0--L8.}
\label{tab:dense_routing}
\small
\begin{tabular}{lccc}
\toprule
Layer & Trained & Untrained & Image\_mean (trained) \\
\midrule
L0 & 10.0\% & 10.0\% & 10.0\% \\
L1 & 100.0\% & 100.0\% & 10.0\% \\
L2 & 100.0\% & 99.6\% & 10.0\% \\
L3 & 100.0\% & 99.6\% & 10.0\% \\
L4 & 100.0\% & 99.1\% & 10.0\% \\
L5 & 100.0\% & 98.2\% & 10.0\% \\
L6 & 100.0\% & 97.3\% & 10.0\% \\
L7 & 100.0\% & 96.4\% & 10.0\% \\
L8 & 100.0\% & 95.5\% & 10.0\% \\
\bottomrule
\end{tabular}
\end{table}

\subsubsection{Cross-Instruction LOOCV Validation}
\label{sec:loocv}

Leave-one-out cross-validation over 10 folds (each fold holds out one task's instructions) confirms the cross-instruction probe result is not driven by memorization.
Trained last\_preaction achieves $100\% \pm 0.00$ across all 10 folds at layers L1--L31, confirming perfect generalization.

\subsubsection{QA Probe: Progress and Subtask Classification}
\label{sec:qa_probe}

To test whether LLM layers encode task-relevant semantic information beyond instruction identity, we train probes to classify task progress (early/mid/late) and active subtask from frozen representations (Table~\ref{tab:qa_probe}).

\begin{table}[h]
\centering
\caption{QA probe accuracy (progress + subtask classification, best layer).}
\label{tab:qa_probe}
\small
\begin{tabular}{lcc}
\toprule
Source & Best Layer & Accuracy \\
\midrule
Trained last\_preaction & L20 & 86.3\% \\
Untrained last\_preaction & L20 & 85.3\% \\
DINOv2 cls\_token & --- & 82.7\% \\
\bottomrule
\end{tabular}
\end{table}

\subsubsection{L1 Attention Pattern Analysis}
\label{sec:attention_pattern}

To understand how training creates robust instruction routing, we analyze per-head attention distributions at L1 (Figure~\ref{fig:attention_heatmap}).

In the trained model, 25 out of 32 attention heads are instruction-dominant: on average they allocate $58\%$ of attention weight to instruction tokens (${\sim}8\%$ of the sequence) and only $20.9\%$ to image tokens at the readout position.
In the untrained model, 0/32 heads are instruction-dominant: attention defaults to image tokens ($87.8\%$) with only $11.7\%$ allocated to instruction tokens.
Training thus learns deliberate instruction routing, converting L1 into a dedicated routing layer.

Both models achieve near-perfect classification at L1 because instruction embeddings from the frozen tokenizer are inherently highly distinguishable: even the untrained model's $11.7\%$ incidental attention provides enough signal for a linear probe.
The critical difference is robustness: trained deliberate routing (58\% attention, 25/32 heads) creates a signal strong enough to survive 30 subsequent layers ($100\%$ at L31), while untrained incidental routing ($11.7\%$ attention, 0/32 dedicated heads) produces a fragile signal that degrades to $82\%$ by L31.

\begin{figure}[h]
\centering
\includegraphics[width=\columnwidth]{figures/fig8_attention.pdf}
\caption{L1 attention from the readout token, per head (sorted by instruction attention fraction). \textbf{(a)}~Trained: 25/32 heads are instruction-dominant ($>50\%$ attention to instruction tokens). \textbf{(b)}~Untrained: 0/32 heads are instruction-dominant; nearly all attention goes to image tokens. Training learns deliberate instruction routing, explaining why the trained signal persists to L31 ($100\%$) while the untrained signal decays ($82\%$).}
\label{fig:attention_heatmap}
\end{figure}
